% HD-F1-I01. Complete proof, independently reviewed at the recorded hash.
\section[HD-F1-I01: the nine-letter four-value rank]{HD-F1-I01: every rational nine-letter four-value component has period rank one}
\label{sec:f1-nine-letter-four-value-rank}

Retain the actual monomial tower, polarization and primitive
character of HD-I01--HD-I21. On every component with
\[
 n=5,\qquad m=9,\qquad g(Y)=0,\qquad b=\nu=4,
 \qquad d\in\{3,4,6\},\quad e\mid d,
\]
the period differential has rank one at every smooth Hurwitz
point. In particular its image has dimension one. This
evaluates the remaining modified-trace-bundle gate for this
entire numerical stratum, including its unramified uncolored
cycles. The stratum is nonempty for $G(3,1,9)$, as the actual
generating tuple below shows.

\paragraph{1. The actual evaluation isomorphism.}
Let $Q$ be the reduced divisor of all uncolored cycles over
the four branch values, exactly as in HD-I20--HD-I21.
The original rank-one Chevalley--Weil calculation gives
$\deg R=2(n-g+1)=12$. Riemann--Hurwitz for the actual
degree-nine map $h:Y\to\mathbf P^1$ gives total
ramification $2m-2=16$. Thus the four fibres have a total
of $4m-16=20$ distinct points, of which twelve are colored.
Consequently $\deg Q=8$. This also agrees with
HD-I20's exact identity
$\deg Q=(\nu-2)m+A-2n=8$, where $A=b-\nu=0$.

For every scheme-theoretic fibre $D=h^*(z)$, of length nine,
the actual restriction sequence is
\[
 0\longrightarrow\mathcal O_Y(Q-D)
  \longrightarrow\mathcal O_Y(Q)
  \longrightarrow\mathcal O_D(Q)\longrightarrow0.
\]
On the genus-zero curve $Y$, its first line bundle has
degree $-1$, so both its $H^0$ and $H^1$ vanish. Restriction
is therefore an isomorphism
$H^0(Y,\mathcal O_Y(Q))\simeq H^0(D,\mathcal O_D(Q))$.
The source has dimension nine. Finite flat pushforward and
the evaluation map consequently give the actual bundle
isomorphism
\[
 H^0(Y,\mathcal O_Y(Q))\otimes\mathcal O_{\mathbf P^1}
       \xrightarrow{\ \sim\ }h_*\mathcal O_Y(Q).
\]
Indeed both sides are locally free of rank nine and the
displayed map is an isomorphism on every fibre; Nakayama's
lemma applies to its kernel and cokernel. Ramification and
nonreduced fibres have been retained in the restriction
sequence.

The literal constant function $1$ is a nonzero vector in
this space of global sections. Extend it to a basis. Under
the actual evaluation isomorphism the constant inclusion
$\mathcal O\to h_*\mathcal O_Y(Q)$ is its first basis
vector. Its quotient is therefore
\[
 \mathcal W_Q=h_*\mathcal O_Y(Q)/\mathcal O
                   \simeq\mathcal O_{\mathbf P^1}^{\oplus8}.
\]
This supplies an actual splitting, including the prescribed
unit section. It follows from HD-I21's actual tangent-fibre
isomorphism that
\[
 \ker d\mathcal F
   \simeq H^0(\mathbf P^1,\mathcal W_Q(2-\nu))
   =H^0(\mathbf P^1,\mathcal O(-2)^{\oplus8})=0.
\]
The Hurwitz component has dimension $b-3=1$, so
$\operatorname{rank}d\mathcal F=1$. The two actual primitive
eigendifferential line bundles have degree four on $Y$,
so their multiplication in HD-I12 is the map
\[
 H^0(\mathcal O_{\mathbf P^1}(4))\otimes
 H^0(\mathcal O_{\mathbf P^1}(4))
       \longrightarrow H^0(\mathcal O_{\mathbf P^1}(8)).
\]
For each $0\le k\le8$, choose integers $a,b$ with
$0\le a,b\le4$ and $a+b=k$. The degree-eight monomial
$X^{8-k}Y^k$ is the product of the degree-four monomials
$X^{4-a}Y^a$ and $X^{4-b}Y^b$. Thus this actual
multiplication is onto at every point. The exact polarized period factorization
in HD-I20--HD-I21 therefore gives
$\operatorname{rank}d\mathrm{Per}=1$ at every such point.

\paragraph{2. A full monomial tuple in this stratum.}
Use colors in $\mathbf Z/3$ for $G(3,1,9)$.
Let $C$ have permutation $(1\,2\,3\,4\,5\,6)$,
cycle color zero on that six-cycle, and colors $1,1,1$
on the fixed letters $7,8,9$. Let $D$ have permutation
$(6\,7\,8\,9)$, cycle color one on that four-cycle,
colors one and two on fixed letters $1,2$, and zero on
fixed letters $3,4,5$. A cycle is represented by one
matrix entry with its stated color and all other entries
equal to one. These specify actual monomial matrices.
Take the original product-one tuple
\[
                    (C,C^{-1},D,D^{-1}).
\]
All four entries are nonreflections.

Here is a complete generation check. Write $c,d$ for the
two underlying permutations. Their commutator is
$(6\,1\,7)$. Conjugating by the powers of $c$ supplies
the triples $(7\,i\,i+1)$ around the six-cycle. These
generate $A_{\{1,\ldots,7\}}$: starting with one
three-cycle, successively add the next letter by a
three-cycle sharing two old letters. To verify this
elementary step, two such three-cycles generate the
alternating group on their four-letter union by
conjugation; thereafter the alternating group on at
least four old letters is transitive on ordered pairs,
so conjugation gives every triple containing the new
letter. Together with the old triples these generate
the alternating group on the enlarged set.
Conjugation by $d$ and then $d^2$ extends this group to
$A_8$ and $A_9$ by the same argument. Since $c$ is odd,
$\langle c,d\rangle=S_9$.

The actual diagonal matrix $D^4$ has color one on
$1,6,7,8,9$, color two on $2$, and zero on $3,4,5$.
Conjugation of diagonal matrices by any monomial matrix
permutes their coordinates. The established surjection
of $\langle C,D\rangle$ to $S_9$ thus supplies a lift
of a transposition exchanging a color-one coordinate
and a color-zero coordinate of $D^4$. Their actual
commutator is a unit coordinate difference, up to sign.
Conjugating it gives all such differences, which span
the total-color-zero subgroup of $(\mathbf Z/3)^9$.
The sum of the colors of $D^4$ is $5+2=1$ modulo three.
It and the differences therefore generate the entire
diagonal group. Canceling the diagonal part of each
permutation lift now gives every ordinary permutation
matrix. Hence $\langle C,D\rangle=G(3,1,9)$.

The two six-cycle permutations contribute ramification
$2\cdot5$, and the two four-cycle permutations contribute
$2\cdot3$. Thus the actual permutation quotient has genus
zero. Each of $C,D$ has exactly three colored cycles, and
each inverse has the three complementary colors. Hence
the actual rank-one cyclic cover of $Y$ has twelve
colored points. Each inverse pair contributes one to the
sum of fractional colors in either primitive eigenspace;
there are six pairs. Chevalley--Weil gives both dimensions
$g-1+6=5$. The Riemann existence correspondence applied
to this product-one full-group tuple gives an actual
connected cover and its monomial quotient tower. This
establishes nonemptiness of the stated rank-one stratum.

\paragraph{3. Retained geometric maps and exact scope.}
The period comparison is the original polarized isogeny
of HD-I02 and the actual primitive multiplication followed
by the canonical quadratic trace of HD-I20. The new map
is the evaluation isomorphism above, followed by quotient
by its specified constant section. The original cover,
finite character, corner and polarization remain the
objects defining that period map. The source is a genuine
one-dimensional family; no dominance in the
twenty-five-dimensional Weil domain is asserted.
Its Hilbert and cycle realizations, Poincare comparison,
and coherent $K_0$/Adams identities are precisely those
of the retained actual tower under F01--F04. This computes
the specified component gate in T4-INDIVIDUAL and leaves
all other numerical strata in their previously stated
exact modified-trace-bundle gate.

\paragraph{Published comparison through the same forgetful map.}
Filom, \emph{On the Rank of the Thurston Pullback Map},
IMRN 2023, no.~24, 21838--21873,
doi:10.1093/imrn/rnac305, Theorem~1.3 (first assertion),
also supplies nonvanishing here. The exact comparison
is as follows. Mark the twelve actual colored points
by a set $A$ and the four branch values by $B$.
The underlying degree-nine oriented branched cover is
an admissible map $(S^2,A)\to(S^2,B)$ in Filom's
definition: $h(A)\subset B$ and $B$ contains all its
critical values. Its marked Hurwitz space has the
actual source-marking map $\rho_A$ and the branch map
$\rho_B$ in his diagram~(2.2). Locally the latter is
an unramified coordinate map. Pulling back the target
complex structure along this same $h$ gives the same
marked source as $\rho_A$; the diagram identifies its
differential with the Thurston pullback differential.
For the original cyclic-data forgetful map, the color
labels are discrete, and the required root line bundle
on $\mathbf P^1$ is unique of its fixed degree.
Its root isomorphism differs only by a scalar, removed
over $\mathbf C$ by rescaling that line bundle.
Consequently the cyclic-data marking adds no tangent
direction to $\rho_A$. These are the actual local
maps defining $d\mathcal F$ above.

Here $|A|=12>9+2$. Filom's first assertion gives
nonzero differential everywhere, hence rank one on
this one-dimensional Hurwitz source. The complete
first-part proof in his Section~3.2 was checked:
the quadratic-trace equations at a regular fibre
give a full-rank Cauchy matrix. Thus nonvanishing
has this published precedent. The explicit evaluation
isomorphism, prescribed-unit quotient, original
polarized period comparison and generating tuple
above give the complete calculation in the project's
retained objects and notation. No priority claim for
the nonvanishing assertion is made.
